%% kapa: exam=1 round=29 year=2018 period=1 form=A subject=real_estate subject_ko=부동산학원론 seq=16 number=96
%% source: Q-Net 감정평가사 기출문제 (cst00309) artlSeq=5202808 file=extracted/1008/5202808/2018년 제29회 감정평가사 1차시험 기출문제/2018년 제29회 감정평가사 1차시험 기출문제/1교시 A형.pdf
%% conversion: kordoc 4.24.51 -> tools/gen1.py (auto); review=proofread-sonnet; flags=none
\begin{question}{96}
우리나라의 부동산투자회사제도에 관한 설명으로 옳지 않은 것은?
\begin{choices}
  \item 자기관리 부동산투자회사의 설립 자본금은 5억원 이상이다.
  \item 부동산투자회사는 발기설립의 방법으로 하여야 하며, 현물출자에 의한 설립이 가능하다.
  \item 위탁관리 부동산투자회사는 자산의 투자ㆍ운용업무를 자산관리회사에 위탁하여야 한다.
  \item 부동산투자회사는 최저자본금준비기간이 끝난 후에는 매 분기 말 현재 총자산의 100분의 80 이상을 부동산, 부동산 관련 증권 및 현금으로 구성하여야 한다.
  \item 부동산투자회사의 상근 임원은 다른 회사의 상근 임직원이 되거나 다른 사업을 하여서는 아니 된다.
\end{choices}
\end{question}
\answer{2}
